\chapter{Detailed EasyCrypt Proof-Code Analysis}
\label{app:easycrypt-code-analysis}

\providecommand{\ecfile}[1]{\nolinkurl{#1}}
\providecommand{\ecsym}[1]{\nolinkurl{#1}}

This appendix records a file-by-file analysis of the EasyCrypt code that is part
of the provable-security verification surface.  The source of truth for the
analysis is the manifest \ecfile{provable-security/proof-files.txt}; files not
listed there are outside the provable-security gate documented here.  In
particular, known-answer-test files for concrete FIPS 202 outputs are useful for
implementation validation but are not needed for the game-based provable-security
argument.

\section{Manifest-level proof-code inventory}
\label{app:ec-inventory}

The manifest contains 16 EasyCrypt files and 22046 lines of checked source.  The
recorded inventory is reproducible with
\ecfile{dissertation/etc/count-easycrypt-declarations.sh}, which fixes the
counting policy used for this appendix and reports axioms separately from
lemmas.  Under the inventory policy used for this dissertation, the scan reports
1088 lemma declarations, 0 axiom declarations, 105 Hoare/equivalence proof
commands, 113 module or module-type declarations, 678 operation/predicate
declarations, and 63 type declarations.  These
numbers are a command-oriented inventory, not a semantic theorem count.  Broader
catalogs that include clones, generated headers, or different Hoare/equiv match
rules may legitimately report different totals, but they should cite their
counting policy explicitly.

The same scan found no literal occurrence of \ecsym{axiom}, \ecsym{admit.}, or
\ecsym{sorry.} in the copied provable-security proof surface.  This does not
remove the cryptographic assumptions from the theorem; it means that those
assumptions are represented as abstract games, hardness terms, modules, and
parameters rather than as unchecked EasyCrypt axioms inside the manifest.

\begingroup
\scriptsize
\setlength{\tabcolsep}{2pt}
\begin{longtable}{p{0.28\textwidth}rrrrrrr}
\caption{Structural summary of the manifest EasyCrypt files.}
\label{tab:easycrypt-code-inventory}\\
\toprule
File & Lines & Types & Ops & Modules & Lemmas & Axioms & Hoare/equiv \\
\midrule
\endfirsthead
\toprule
File & Lines & Types & Ops & Modules & Lemmas & Axioms & Hoare/equiv \\
\midrule
\endhead
\ecfile{HAETAE\_Security.ec} & 1272 & 0 & 1 & 0 & 40 & 0 & 0 \\
\ecfile{Sig\_ROM.ec} & 206 & 11 & 4 & 16 & 3 & 0 & 0 \\
\ecfile{HAETAE\_HopGames.ec} & 8817 & 1 & 40 & 62 & 247 & 0 & 88 \\
\ecfile{HAETAE\_Reductions.ec} & 504 & 0 & 26 & 0 & 50 & 0 & 0 \\
\ecfile{HAETAE\_ROM.ec} & 111 & 4 & 24 & 2 & 3 & 0 & 0 \\
\ecfile{HAETAE\_ROM\_Programming.ec} & 1754 & 2 & 43 & 11 & 99 & 0 & 17 \\
\ecfile{HAETAE\_Scheme.ec} & 66 & 7 & 2 & 1 & 0 & 0 & 0 \\
\ecfile{HAETAE\_Transcript.ec} & 660 & 1 & 51 & 0 & 35 & 0 & 0 \\
\ecfile{HAETAE\_Events.ec} & 57 & 0 & 4 & 0 & 4 & 0 & 0 \\
\ecfile{HAETAE\_Assumptions.ec} & 377 & 5 & 13 & 21 & 11 & 0 & 0 \\
\ecfile{HAETAE\_Distributions.ec} & 1218 & 2 & 64 & 0 & 97 & 0 & 0 \\
\ecfile{HAETAE\_Rejection.ec} & 2204 & 4 & 97 & 0 & 139 & 0 & 0 \\
\ecfile{HAETAE\_Algebra.ec} & 3885 & 20 & 246 & 0 & 282 & 0 & 0 \\
\ecfile{HAETAE\_Params.ec} & 113 & 1 & 17 & 0 & 19 & 0 & 0 \\
\ecfile{HAETAE\_FIPS202.ec} & 105 & 2 & 12 & 0 & 9 & 0 & 0 \\
\ecfile{HAETAE\_Keccak1600.ec} & 697 & 3 & 34 & 0 & 50 & 0 & 0 \\
\midrule
Total & 22046 & 63 & 678 & 113 & 1088 & 0 & 105 \\
\bottomrule
\end{longtable}
\endgroup

The proof surface has four strata.  The foundation stratum defines parameters,
algebraic data structures, SHAKE/Keccak support, and distributions.  The scheme
stratum defines the HAETAE random-oracle scheme interface and the signing,
verification, transcript, event, and rejection predicates.  The game stratum
defines EUF-CMA, UF-NMA, ROM, transcript, sampler, and programming games.  The
reduction stratum composes hop bounds into the final checked implication.

\section{Trust boundary and checked claims}
\label{app:ec-trust-boundary}

The machine-checked claim is a conditional game-based security implication.  The
checked proof establishes that the fixed-ledger EasyCrypt-model EUF-CMA
advantage in the modeled random-oracle setting is bounded by the concrete
expression assembled in \ecfile{HAETAE\_Reductions.ec}, subject to the explicit
final hop/reduction premise, the abstract hardness games, and the parameter
assumptions represented in \ecfile{HAETAE\_Assumptions.ec}.  One current
final-hop premise is a
proof-surface premise rather than a non-vacuous concrete-security premise because
it contains the oversized \ecsym{rejection\_sampling\_loss\_term}.  Thus the
checker verifies the game transformations, sampler/event accounting, algebraic
well-formedness obligations, and non-negativity/subunit arithmetic needed to
compose the final implication.  It does not prove Module-SIS hardness, MLWE
hardness, the missing non-vacuous final hop, or the physical correctness of an
implementation outside the modeled operations.

The absence of manifest-local \ecsym{axiom}, \ecsym{admit.}, and
\ecsym{sorry.} occurrences is important because the final implication is not
obtained by installing an unchecked theorem statement.  The remaining trusted
base consists of EasyCrypt's logic and libraries, the correctness of the
manifest selection, the intended interpretation of abstract hardness games, and
the fidelity of the HAETAE model to the specification.

\section{File-by-file analysis}
\label{app:ec-file-analysis}

\subsection[HAETAE\_Security.ec]{\texorpdfstring{\ecfile{HAETAE\_Security.ec}}{HAETAE\_Security.ec}}

\ecfile{HAETAE\_Security.ec} is the top-level composition file.  It imports the
signature ROM interface, the HAETAE scheme model, algebra, distributions,
reductions, assumptions, rejection logic, transcript logic, ROM programming,
and hop games.  Its purpose is compositional: it does not define a new
low-level model, but connects previously checked lemmas into externally usable
security statements.

The first obligation is correctness.  The operation
\ecsym{correctness\_obligation} packages the required rejection-freeness and
verification facts, while \ecsym{correctness\_obligation\_holds} and
\ecsym{correctness\_perfect} connect those facts to the correctness game.  This
establishes that the modeled signing algorithm produces signatures accepted by
the modeled verifier under the proof assumptions.

The EUF-CMA part is organized as a chain of increasingly explicit theorems.
Lemmas such as \ecsym{euf\_cma\_security\_no\_signing},
\ecsym{euf\_cma\_security}, and
\ecsym{euf\_cma\_security\_with\_hop\_interfaces} expose the public theorem at
different levels of abstraction.  The middle of the file refines the proof
through simulated-signing, explicit-boundary, ROM-internal transcript, clear
site, counted clear site, structural NMA, public simulation, paper simulation,
real-signing paper simulation, HAETAE-rejection paper simulation, exact
hyperball, and RO-exact hyperball stages.  The final lemmas, including
\ecsym{euf\_cma\_security\_from\_paper\_sample\_lifting\_and\_counted\_bimodal},
\ecsym{euf\_cma\_security\_from\_budgeted\_paper\_sample\_lifting\_and\_counted\_bimodal},
\ecsym{euf\_cma\_security\_from\_checked\_ro\_sampling\_hop\_and\_counted\_bimodal},
and
\ecsym{euf\_cma\_security\_from\_real\_signing\_exact\_hyperball\_hop\_and\_counted\_bimodal},
are the checked composition points used by the dissertation security
implication.

The code pattern in this file is deliberately conservative.  Most lemmas are
bound-lifting or theorem-composition lemmas whose proof obligations are
arithmetic inequalities and applications of earlier hop lemmas.  This makes the
top-level implication auditable: each public statement records which lower-level
boundary it depends on, and the concrete loss expression remains centralized in
\ecfile{HAETAE\_Reductions.ec}.

\subsection[Sig\_ROM.ec]{\texorpdfstring{\ecfile{Sig\_ROM.ec}}{Sig\_ROM.ec}}

\ecfile{Sig\_ROM.ec} defines the generic signature security interface used by
the HAETAE proof.  It abstracts public keys, secret keys, messages, contexts,
signatures, random-oracle queries, and random-oracle outputs.  The operations
\ecsym{fresh\_msg} and \ecsym{fresh\_sig} define freshness predicates for
EUF-CMA and strong-EUF-CMA style games, and the file clones EasyCrypt's full
random-oracle interface through \ecsym{FullRO}.

The module types \ecsym{Oracle}, \ecsym{POracle}, \ecsym{Scheme},
\ecsym{CORR\_ADV}, \ecsym{SignOracle}, \ecsym{Adversary}, and
\ecsym{NMA\_Adversary} are the interfaces through which all later HAETAE games
are parameterized.  The game modules \ecsym{Correctness}, \ecsym{EUF\_CMA},
\ecsym{UF\_NMA}, and \ecsym{SUF\_CMA} define the standard experiments.  The
adapter \ecsym{NMA\_As\_CMA} embeds a no-message adversary in a chosen-message
interface.

The central proof obligation in this file is
\ecsym{nma\_as\_cma\_exact}.  It is an equivalence proof showing that the NMA
adapter preserves the result of the experiment under equal globals.  Later
HAETAE reductions rely on this generic lemma when replacing active chosen-message
interfaces with passive or simulated transcript interfaces.

\subsection[HAETAE\_HopGames.ec]{\texorpdfstring{\ecfile{HAETAE\_HopGames.ec}}{HAETAE\_HopGames.ec}}

\ecfile{HAETAE\_HopGames.ec} is the largest and most security-critical proof
file in the manifest.  It contains the detailed game transformations that turn
real HAETAE signing into transcript-based and ROM-internal simulators, then into
NMA samplers, paper samplers, rejection samplers, exact-hyperball samplers, and
budgeted random-oracle lifting games.

The first block defines signing simulators and transcript consistency
invariants.  Predicates such as \ecsym{transcript\_log\_consistent},
\ecsym{transcript\_log\_ready}, \ecsym{internal\_transcript\_state\_sound}, and
\ecsym{sampler\_rom\_covered} are the invariants that let the proof state that
logged transcripts match signatures, public keys, mode data, min-entropy
conditions, and random-oracle query coverage.  Lemmas including
\ecsym{transcript\_log\_consistent\_cons},
\ecsym{internal\_transcript\_state\_sound\_sign\_internal\_cons}, and
\ecsym{sampler\_rom\_covered\_fresh\_after\_clean\_seed} are preservation
lemmas used in Hoare proofs for signing oracles.

The next major block defines game modules for simulated signing:
\ecsym{EUF\_CMA\_SimulatedSign},
\ecsym{EUF\_CMA\_TranscriptSimulatedSign},
\ecsym{EUF\_CMA\_InternalTranscriptSign}, and
\ecsym{EUF\_CMA\_ROMInternalTranscriptSign}.  These modules isolate the
semantic changes between real signing and transcript-oriented simulation.  The
proof obligations here are mainly equivalence and Hoare obligations showing
that results and bad events are preserved or bounded when the signing oracle is
replaced.

The paper-simulation block defines simulated transcript and signature samples,
including public-simulation fields and RO-signing-attempt sampling.  Operations
such as \ecsym{paper\_sim\_sample\_token},
\ecsym{public\_sim\_signature}, and the RO signing-attempt samplers encode the
paper proof's distributional objects in a form that EasyCrypt can reason about.
This is the point where the proof maps hand-written sampling arguments into
typed programs and event predicates.

The budgeted ROM-internal block is the most delicate part of the file.  It
defines budgeted wrappers for adversary hash calls and signing calls, logs
adversary-side hash queries and sampler self-queries, and introduces
sampler-clean events such as \ecsym{sampler\_bad\_prequery}.  The checked
lifting theorems compare
\ecsym{ROSigningAttemptPaperSimSampler} and
\ecsym{ROExactHyperballPaperSimSampler} under clean conditions and query
budgets.  The important proof idea is a fundamental-lemma argument: the attempt
side and exact side are equivalent unless the sampler's fresh
\ecsym{sampler\_expand\_query} was already present in the joint adversary/self
log.

The concentration of 88 Hoare/equivalence commands in this file is expected.
Those commands correspond to the stateful game transitions that cannot be
captured by pure algebraic rewriting alone.  This file is therefore the primary
machine-checked replacement for the informal hybrid-game portion of the
pen-and-paper proof.

\subsection[HAETAE\_Reductions.ec]{\texorpdfstring{\ecfile{HAETAE\_Reductions.ec}}{HAETAE\_Reductions.ec}}

\ecfile{HAETAE\_Reductions.ec} defines the concrete loss terms that appear in
the final checked implication.  It introduces abstract hardness terms such as
\ecsym{mlwe\_hardness\_term}, \ecsym{module\_sis\_hardness\_term}, and
\ecsym{bimodal\_selftarget\_msis\_hardness\_term}, then defines query budgets,
challenge-support lower bounds, counted ROM collision/prequery/min-entropy
terms, rejection-sampling terms, and final EUF bounds.

The code is mostly arithmetic normalization and inequality proof.  Lemmas such
as \ecsym{haetae\_euf\_boundE},
\ecsym{haetae\_euf\_bound\_groupedE},
\ecsym{haetae\_euf\_counted\_rom\_bound\_groupedE},
\ecsym{counted\_rom\_programming\_loss\_term\_subunit},
\ecsym{counted\_rom\_support\_dominates\_budget}, and
\ecsym{counted\_rom\_loss\_from\_query\_budgets\_and\_support} ensure that the
bound is algebraically well-formed, nonnegative, and expressible in terms of the
declared budgets and support sizes.

This file is the numerical interface between proof and theorem statement.  The
proof scripts in the game files produce losses in structured form, while this
file rewrites them into the concrete HAETAE bound used by
\ecfile{HAETAE\_Security.ec}.

\subsection[HAETAE\_ROM.ec]{\texorpdfstring{\ecfile{HAETAE\_ROM.ec}}{HAETAE\_ROM.ec}}

\ecfile{HAETAE\_ROM.ec} specializes EasyCrypt's full random oracle to HAETAE.
It defines the query tags for message hashing, challenge hashing, matrix
expansion, and sampler expansion.  Its random-oracle output type contains a CRH,
a challenge, a matrix, and signing randomness, with projection operations such
as \ecsym{ro\_output\_to\_crh},
\ecsym{ro\_output\_to\_challenge}, \ecsym{ro\_output\_to\_matrix}, and
\ecsym{ro\_output\_to\_random\_coins}.

The distribution \ecsym{dro\_output} is query-dependent: matrix expansion uses
the matrix distribution, sampler expansion uses signing-randomness output, and
other query forms use the appropriate projections.  The lemmas
\ecsym{matrix\_query\_distribution\_lossless},
\ecsym{ro\_output\_distribution\_lossless}, and
\ecsym{sampler\_expand\_query\_dro\_output\_ro\_signing\_coins} establish the
losslessness and projection behavior that the sampler-freshness proof needs.

This file is small but foundational.  If query constructors were not separated,
the later prequery analysis could not state precisely that only adversary
queries equal to \ecsym{sampler\_expand\_query} collide with the signer's fresh
sampler seed.

\subsection[HAETAE\_ROM\_Programming.ec]{\texorpdfstring{\ecfile{HAETAE\_ROM\_Programming.ec}}{HAETAE\_ROM\_Programming.ec}}

\ecfile{HAETAE\_ROM\_Programming.ec} defines counted ROM programming.  It
introduces programming sites, programming-site queries and outputs, transcript
programming sites, challenge entropy checks, min-entropy failures, and counted
ROM events.  These definitions allow the proof to count hash queries, signature
sites, programmed sites, and min-entropy checks in a single event log.

The module \ecsym{CountedLazyROM} is the concrete lazy random oracle used by the
programming proof.  Preservation lemmas such as
\ecsym{counted\_lazy\_rom\_init\_clears},
\ecsym{counted\_lazy\_rom\_get\_preserves\_clear},
\ecsym{counted\_lazy\_rom\_get\_preserves\_bad\_clear},
\ecsym{counted\_lazy\_rom\_observe\_transcript\_preserves\_clear}, and
\ecsym{counted\_lazy\_rom\_program\_preserves\_bad\_clear\_if\_fresh} are the
state invariants used to justify ROM programming without double-programming or
programming a prequeried site.

The later lemmas in the file turn event-counting disciplines into probability
bounds.  They are the bridge between operational ROM code and the arithmetic
loss terms in \ecfile{HAETAE\_Reductions.ec}.  The file therefore supports both
the classical Fiat--Shamir with aborts programming step and the separate
budgeted sampler prequery analysis used in the active-signing lifting surface.

\subsection[HAETAE\_Scheme.ec]{\texorpdfstring{\ecfile{HAETAE\_Scheme.ec}}{HAETAE\_Scheme.ec}}

\ecfile{HAETAE\_Scheme.ec} instantiates the generic signature interface with
HAETAE types.  It clones \ecfile{Sig\_ROM.ec} with HAETAE public keys, secret
keys, messages, contexts, signatures, ROM queries, and ROM outputs.  The
operation \ecsym{haetae\_mode} leaves the parameter mode abstract so that the
security proof applies uniformly across the supported HAETAE parameter sets.

The module \ecsym{HAETAE} supplies the key-generation, signing, and verification
procedures used by the generic games.  Its role is interface alignment, not
cryptographic reasoning.  The proofs about the behavior of signing and
verification are discharged in the algebra, rejection, transcript, and security
files that import this scheme module.

\subsection[HAETAE\_Transcript.ec]{\texorpdfstring{\ecfile{HAETAE\_Transcript.ec}}{HAETAE\_Transcript.ec}}

\ecfile{HAETAE\_Transcript.ec} defines the transcript abstraction used by the
forking, ROM-programming, and rejection-sampling arguments.  Transcript fields
record the public key, message, context, commitment data, challenge, response,
hint, signature token, rejection branch, and validity conditions needed to relate
a signature to the algebraic equations seen by the verifier.

The lemmas in this file prove that transcript predicates expose the exact fields
needed by later game hops.  Examples include
\ecsym{transcript\_fields\_match\_signature\_challenge\_query},
\ecsym{transcript\_valid\_signature\_challenge\_query},
\ecsym{transcript\_valid\_public\_equation},
\ecsym{transcript\_valid\_commitment\_reconstruction}, and
\ecsym{valid\_forking\_pair\_public\_equations}.  The extraction lemmas
\ecsym{extract\_bimodal\_forking\_some} and
\ecsym{extract\_bimodal\_forking\_nonempty} connect valid transcript pairs to
the bimodal self-target MSIS reduction interface.

This file is the semantic link between game-based forgery events and algebraic
hardness instances.  Without transcript validity lemmas, a final forgery would
remain an accepted signature only; with these lemmas, the proof can extract the
structured algebraic relation required by the reduction.

\subsection[HAETAE\_Events.ec]{\texorpdfstring{\ecfile{HAETAE\_Events.ec}}{HAETAE\_Events.ec}}

\ecfile{HAETAE\_Events.ec} defines the small collection of event predicates used
to state rejection-freeness and acceptance properties.  The operations describe
key-generation rejection, signing rejection, verification rejection, and accepted
transcript rejection-free conditions.

The four lemmas in the file establish that the current model's key-generation,
signing, verification, and accepted transcript paths satisfy the corresponding
event predicates.  This isolates event vocabulary from the larger proof files:
game hops can refer to named events rather than re-expanding rejection and
verification definitions at every boundary.

\subsection[HAETAE\_Assumptions.ec]{\texorpdfstring{\ecfile{HAETAE\_Assumptions.ec}}{HAETAE\_Assumptions.ec}}

\ecfile{HAETAE\_Assumptions.ec} defines the abstract hardness games and
reduction interfaces used by the final checked implication.  It contains module types for
solvers and reductions, games for Module-SIS, bimodal self-target MSIS, MLWE,
and distribution-lifted relations, plus adapter modules that lift one solver
interface into another.

The lemmas are exactness and well-formedness results for these adapters.  For
example, \ecsym{bimodal\_solver\_lift\_exact},
\ecsym{bimodal\_mode\_solver\_lifted\_distribution\_exact}, and
\ecsym{bimodal\_msis\_as\_module\_sis\_exact} are equivalence proofs showing
that adapter games preserve success events.  Lemmas such as
\ecsym{module\_sis\_eval\_wf}, \ecsym{module\_sis\_valid\_target\_wf}, and
\ecsym{bimodal\_to\_module\_sis\_valid} provide structural well-formedness for
instances passed to these games.

This file is the explicit cryptographic trust boundary.  The machine checker
verifies that HAETAE's game reductions reduce to these games and terms; it does
not prove that the abstract hardness terms are small.  That interpretation is a
mathematical and cryptographic assumption outside EasyCrypt's deductive core.

\subsection[HAETAE\_Distributions.ec]{\texorpdfstring{\ecfile{HAETAE\_Distributions.ec}}{HAETAE\_Distributions.ec}}

\ecfile{HAETAE\_Distributions.ec} defines all distributions needed for the
scheme model and sampler proofs.  It includes byte, seed, CRH, signing
randomness, coefficient, polynomial, vector, matrix, challenge, structural
signing sample, bounded-unit sample, checked-hyperball sample, and
exact-hyperball sample distributions.

The early lemmas prove losslessness and domain correctness for byte, seed, CRH,
and signing randomness distributions.  The signing-randomness token machinery is
especially important: it gives the proof a finite support with a point bound,
which is later used to bound the probability that a fresh sampler seed collides
with a logged adversary query.  Lemmas such as
\ecsym{signing\_coin\_distribution\_lossless},
\ecsym{signing\_entropy\_token\_to\_coins\_tokenK}, and
\ecsym{signing\_randomness\_from\_source\_tokenK} justify that representation.

The hyperball distribution lemmas support the paper-sample to exact-hyperball
transition.  They prove well-formedness, boundedness, losslessness, and
probability point bounds for the distributions that replace the rejection-based
signing attempt distribution.  These lemmas are consumed by
\ecfile{HAETAE\_Rejection.ec} and \ecfile{HAETAE\_HopGames.ec}.

\subsection[HAETAE\_Rejection.ec]{\texorpdfstring{\ecfile{HAETAE\_Rejection.ec}}{HAETAE\_Rejection.ec}}

\ecfile{HAETAE\_Rejection.ec} formalizes the signing-attempt state, rejection
conditions, simulation predicates, and distributional bridges for HAETAE's
rejection sampling.  The file defines fields extracted from signing attempts,
including sampled vectors, commitment high and low bits, challenge, response,
auxiliary response, hint, signature token, rejection-branch bit, and programmed
challenge data.

The first proof block connects signing attempts to valid signatures.  Lemmas
such as \ecsym{signing\_attempt\_verifies},
\ecsym{signing\_simulation\_verifies}, and
\ecsym{signing\_attempt\_simulates} show that states satisfying the modeled
relations produce signatures accepted by verification.  These facts are used by
the correctness theorem and by the game hops that replace concrete signing with
simulated signing.

The sampler block proves that structural, bounded-unit, checked-hyperball, and
exact-hyperball samples satisfy the predicates needed by the signing-attempt
state.  Lemmas such as
\ecsym{structural\_signing\_sample\_pair\_sampler\_ok},
\ecsym{bounded\_unit\_signing\_sample\_pair\_sampler\_ok},
\ecsym{checked\_hyperball\_signing\_sample\_pair\_sampler\_ok}, and
\ecsym{exact\_hyperball\_signing\_sample\_pair\_sampler\_ok} are the checked
counterparts of the paper proof's sampling claims.

The rejection-bound block connects abort predicates to the global
\ecsym{rejection\_sampling\_loss\_term}.  This is where the proof accounts for
the cost of moving from real rejection sampling to the idealized distributions
used in the later NMA and exact-hyperball games.

\subsection[HAETAE\_Algebra.ec]{\texorpdfstring{\ecfile{HAETAE\_Algebra.ec}}{HAETAE\_Algebra.ec}}

\ecfile{HAETAE\_Algebra.ec} defines the algebraic model for HAETAE.  It gives
types for bytes, seeds, CRHs, random coins, messages, contexts, coefficients,
polynomials, challenges, polynomial vectors, matrices, public keys, secret keys,
encoded public keys, signatures, and signature tokens.  It also defines zero
values, well-formedness predicates, modular coefficient arithmetic, polynomial
arithmetic, vector arithmetic, matrix-vector multiplication, packing-like
operations, high/low-bit operations, challenge operations, and verification
relations.

The 282 lemmas in this file are mostly preservation, size, well-formedness, and
equational lemmas.  Examples include \ecsym{poly\_zero\_wf},
\ecsym{polyveck\_zero\_wf}, \ecsym{poly\_add\_wf},
\ecsym{poly\_mul\_wf}, \ecsym{poly\_dot\_wf}, and the corresponding vector and
matrix facts.  These lemmas are not merely implementation details: later
transcript and rejection proofs depend on them to ensure that every sampled or
derived value has the correct dimension and belongs to the modeled algebraic
domain.

The file is intentionally explicit and list-based.  That representation keeps
the proof close to executable data structures and avoids hiding dimension
conditions in an abstract ring interface.  The cost is a large number of size and
well-formedness obligations, which EasyCrypt checks once and makes reusable
throughout the proof.

\subsection[HAETAE\_Params.ec]{\texorpdfstring{\ecfile{HAETAE\_Params.ec}}{HAETAE\_Params.ec}}

\ecfile{HAETAE\_Params.ec} defines the parameter constants and mode-dependent
functions for HAETAE.  Constants include seed size, CRH size, polynomial degree,
modulus, doubled modulus, and the parameter modes \ecsym{Mode2},
\ecsym{Mode3}, and \ecsym{Mode5}.  Mode functions define vector dimensions,
challenge weight, rejection bounds, hint parameters, signature sizes, and public
key sizes.

The lemmas prove positivity and basic arithmetic facts about these parameters.
Examples include \ecsym{seedbytes\_gt0}, \ecsym{crhbytes\_gt0},
\ecsym{n\_gt0}, \ecsym{q\_gt0}, \ecsym{mode\_k\_gt0},
\ecsym{mode\_l\_gt0}, \ecsym{mode\_tau\_le\_n}, and
\ecsym{mode\_publickeybytes\_gt0}.  These facts discharge side conditions in
distribution losslessness, vector sizing, and probability-bound proofs.

This file is part of the trusted modeling layer because it states the numeric
parameters used by all later code.  EasyCrypt checks consequences of these
definitions, but the correspondence between these constants and the intended
HAETAE parameter sets is a modeling obligation.

\subsection[HAETAE\_FIPS202.ec]{\texorpdfstring{\ecfile{HAETAE\_FIPS202.ec}}{HAETAE\_FIPS202.ec}}

\ecfile{HAETAE\_FIPS202.ec} models the SHAKE256 portion of FIPS 202 used by the
HAETAE algebraic model.  It imports the Keccak permutation model, defines the
state size, SHAKE256 rate, capacity, domain separator, final padding byte,
absorb operation, squeeze operation, and complete SHAKE256 operation.

The lemmas prove constant equalities and output-size properties, including
\ecsym{fips202\_state\_bytesE}, \ecsym{shake256\_rate\_bytesE},
\ecsym{shake256\_capacity\_bytesE},
\ecsym{shake256\_rate\_capacity\_state\_bytesE},
\ecsym{shake256\_absorb\_once\_short\_state\_size},
\ecsym{shake256\_squeeze\_size}, and \ecsym{shake256\_size}.  These facts are
used by the algebraic model to ensure byte-string sizes match the expected
HAETAE encodings.

For provable security, the random oracle abstraction is the main security
object.  This file contributes concrete hash-support definitions used by the
model and by correspondence with implementation-level code, but the game-based
ROM theorem does not rely on known-answer tests.

\subsection[HAETAE\_Keccak1600.ec]{\texorpdfstring{\ecfile{HAETAE\_Keccak1600.ec}}{HAETAE\_Keccak1600.ec}}

\ecfile{HAETAE\_Keccak1600.ec} defines a list-based model of the Keccak-f1600
permutation.  It introduces byte, lane, and state types; normalization and bit
projection operations; lane Boolean operations; rotations; byte-lane
conversions; state indexing; round constants; and the theta, rho-pi, chi, and
iota round functions.

The lemmas prove well-formedness and bit-level equations for these operations.
Examples include \ecsym{keccak\_lane\_zero\_wf},
\ecsym{keccak\_lane\_xor\_wf}, \ecsym{keccak\_lane\_and\_wf},
\ecsym{keccak\_lane\_rotl\_wf}, \ecsym{keccak\_bytes\_to\_lane\_wf},
\ecsym{keccak\_theta\_bitE}, \ecsym{keccak\_rho\_pi\_bitE}, and
\ecsym{keccak\_chi\_bitE}.  These lemmas support the FIPS 202 model's size and
bit-level correctness obligations.

In the provable-security proof, this file is support code rather than the main
cryptographic reduction.  Its function is to make the hash-related HAETAE model
internally consistent.  Security is still proved in the ROM abstraction; the
Keccak model documents and checks the concrete hash representation used by the
formalized scheme layer.

\section{Cross-file dependency analysis}
\label{app:ec-cross-file-analysis}

The dependency flow is acyclic at the conceptual level.  Parameters feed
algebra; algebra and parameters feed distributions; parameters, algebra, and
distributions feed the ROM and scheme; algebra, distributions, events, and
transcripts feed rejection; ROM and transcript logic feed ROM programming; ROM
programming, rejection, transcript logic, and reductions feed hop games; hop
games and reductions feed the top-level security theorem.

The most important design choice is the separation between pure predicates and
stateful games.  Pure files such as \ecfile{HAETAE\_Params.ec},
\ecfile{HAETAE\_Algebra.ec}, \ecfile{HAETAE\_Distributions.ec},
\ecfile{HAETAE\_Transcript.ec}, and \ecfile{HAETAE\_Reductions.ec} define
operations and prove algebraic or probabilistic facts about them.  Stateful files
such as \ecfile{Sig\_ROM.ec}, \ecfile{HAETAE\_ROM\_Programming.ec}, and
\ecfile{HAETAE\_HopGames.ec} define games and use Hoare or equivalence proofs to
relate adversarial executions.  This split is useful because it localizes
program-logic obligations to the files where mutable game state is actually
present.

The second important design choice is explicit event logging.  Transcript logs,
ROM query logs, counted ROM event logs, and sampler self-query logs make the
fundamental-lemma proofs inspectable.  The budgeted sampler prequery event is not
a black-box union bound over an oversized wrapper; it is tied to the concrete
event that a fresh \ecsym{sampler\_expand\_query} is already present in the
logged joint log of adversary hash queries and previous sampler self-queries.
This event is distinct from the \ecsym{counted\_rom\_prequery\_term} summand in
the final counted-ROM bound.

The third important design choice is centralization of concrete arithmetic.
Loss terms are not duplicated across hop proofs.  Hop files prove local
probability inequalities, and \ecfile{HAETAE\_Reductions.ec} normalizes them into
the final bound.  This reduces the risk that two theorem statements silently use
different interpretations of the same loss.

\section{Residual proof obligations outside EasyCrypt}
\label{app:ec-residual-obligations}

The analysis identifies four residual obligations that are outside the manifest
even though the manifest itself is machine checked.  First, the abstract hardness
terms in \ecfile{HAETAE\_Assumptions.ec} must be interpreted as the intended
MLWE, Module-SIS, and bimodal self-target MSIS assumptions.  Second, the
parameter constants in \ecfile{HAETAE\_Params.ec} must match the parameter set
claimed by the surrounding HAETAE specification.  Third, the modeled signing,
verification, rejection, and transcript operations must be kept aligned with the
implementation and the paper specification.  Fourth, the manifest
\ecfile{provable-security/proof-files.txt} must remain the audited proof surface;
adding or removing EasyCrypt files changes the scope of the claim and requires a
fresh inventory and compile check.

Within those boundaries, the EasyCrypt code performs the intended
machine-checked provable-security work: it defines the security games, formalizes
the HAETAE model, proves the transcript and sampler invariants, carries out the
ROM and budgeted lifting hops, proves the concrete arithmetic side conditions,
and composes the checked lemmas into the final conditional EUF-CMA implication.
The
main dissertation now treats this appendix as evidence, not as a substitute for
the mathematical proof narrative; curated proof-object analyses appear in
Appendix~\ref{app:curated-easycrypt-case-studies}, and the exhaustive generated
individual correspondence is archived under \ecfile{dissertation/etc/}.
